\subsection{Overview of AdsMind}

AdsMind implements a closed-loop agent architecture for adsorption configuration search. The system comprises five modules: \textbf{Planner}, \textbf{Validator}, \textbf{Executor}, \textbf{Analyzer}, and \textbf{Summarizer} (Figure~\ref{fig:pipeline}). In each iteration, the \textbf{Planner} proposes an adsorption hypothesis, which the \textbf{Validator} checks before physical execution. The \textbf{Executor} relaxes the structure, and the \textbf{Analyzer} interprets the post-relaxation coordinates to guide the next iteration. The \textbf{Summarizer} routes control flow and emits the final report when termination criteria are met.

\subsection{Closed-Loop Agent Architecture}

\subsubsection{Planner and Validator: Adsorption Hypothesis Generation}

At the beginning of the analysis, the user provides a request prompt together with the corresponding adsorbate SMILES string and surface slab file.
At the onset of each iteration, the \textbf{Planner} proposes a structured adsorption hypothesis in JSON format, including the canonical site type---\textit{ontop} (one surface atom), \textit{bridge} (two adjacent surface atoms), or \textit{hollow} (three or more surface atoms)---the intended surface symbols, and the adsorbate binding indices~\cite{norskov2004density,shriver2006inorganic}.

In our closed-loop architecture, the updated context from other modules is entered into the \textbf{Planner} prompt, providing physical feedback and performance analysis from previous hypotheses.
The \textbf{Planner} prompt is assembled from conditional sections rather than a monolithic, fixed template. Three runtime switches control the agentic feedback exposed to the language model:
\textit{Slip feedback} (enables slip-conditioned guidance) determines whether prior slips are framed as explicit evidence that the originally planned site is unstable.
\textit{Forbid} (explicit negative constraint) determines whether slipped or otherwise failed site intents are converted into an explicit prohibition in the planner history, preventing retry of the same failed site family.
\textit{Termination} determines whether convergence-to-known-best signals are surfaced to the planner and whether the planner is allowed to stop early after repeated convergence.
The three switches are independent; disabling all three yields the 1-Shot baseline. When Termination is disabled, the loop continues until the attempt budget (default: 5) is exhausted.

Convergence is declared when the same site family is proposed in three consecutive iterations without finding a lower-energy configuration. The known-best energy is the minimum adsorption energy observed across all iterations. When Termination is enabled, the Planner receives a summary of the best-so-far result and may emit a STOP signal.

%% Validator
After the hypothesis is proposed by the \textbf{Planner}, the \textbf{Validator} enforces chemistry- and schema-level constraints before any expensive physics simulation is made. The validation protocol ensures that:
\begin{enumerate}
    \item The requested adsorbate binding indices are valid and correspond to non-hydrogen heavy atoms, unless the adsorbate is explicitly a single atom (e.g., [H]).
    \item The cardinality of the binding indices matches the topological requirements of the specified site type (e.g., \textit{ontop} requires exactly one index, while \textit{bridge} allows up to two indices for side-on modes).
    \item The hypothesis has not been previously evaluated within the current search trajectory, thereby preventing redundant calculations.
\end{enumerate}
Hypotheses failing these checks elicit an immediate penalty feedback prompt, directing the LLM to refine its proposal without advancing the state machine to the physical simulation phase.

\subsubsection{Executor and Analyzer: Physical Simulation and Scientific Analysis}

%% Executor
With the validated adsorption hypothesis, the \textbf{Executor} performs physical simulation to evaluate the adsorption energy and performance.
In AdsMind, the \textbf{Executor} instantiates the accepted plan on the slab, performs physical relaxation, and returns a structured post-relaxation analysis rather than only a scalar energy.
The \textbf{Analyzer} records whether the final structure remained bound, whether dissociation or rearrangement occurred, and whether the realized coordination environment differs from the \textbf{Planner}'s intent. These deviations are logged as slip diagnostics. \textit{Chemical Slip} is detected when the relaxed structure exhibits a different surface-symbol fingerprint than the Planner's intent. The surface-symbol fingerprint is defined as the sorted list of element symbols of surface atoms within the first coordination shell of the adsorbate anchor atom. Canonical site-family mismatch is retained separately as a geometric/site slip signal in the search history.

Before the physical simulation, the proposal is pre-processed by being translated into an atomistic starting structure through a surrogate-SMILES geometry layer. 
This layer maps the selected binding atom(s) to proxy adsorption markers and constructs physically valid initial placements for \textit{ontop}, \textit{bridge}, and \textit{hollow} adsorption.
It samples a capped set of candidate surface sites and conformers.
In the benchmark configuration, the default cap is eight filtered surface coordinates (matching the requested site type and surface-symbol pattern) and four conformers per site, with one top-ranked candidate carried into full relaxation.

A \textbf{two-step evaluation procedure} is then employed.
Candidate structures are first screened with the MACE-MP calculator used by the active protocol to discard pathological configurations and rank promising candidates by screened energy. This screening stage can optionally include a short Langevin molecular-dynamics warmup, but the warmup is disabled (\texttt{md\_steps=0}) in the primary CPU benchmark. The best surviving candidate (\texttt{relax\_top\_n=1}) is then fully relaxed and analyzed.
The post-relaxation analysis reconstructs the actual local coordination environment, detects dissociation and bond changes, and writes structured artifacts that can be consumed directly by the next \textbf{Planner} step or by offline benchmark tooling.

For a given adsorption state, the thermodynamic stability is quantified by the adsorption energy $E_{\text{ads}}$, defined as:
\begin{equation}
    E_{\text{ads}} = E_{\text{total}} - E_{\text{surface}} - E_{\text{adsorbate}}
    \label{eq:adsorption_energy}
\end{equation}
where $E_{\text{total}}$ is the total potential energy of the relaxed adsorbate-slab system, $E_{\text{surface}}$ is the single-point potential energy of the bare surface evaluated with the same cell and calculator settings, and $E_{\text{adsorbate}}$ is the energy of the isolated adsorbate in a vacuum box.

Structural relaxation is performed using the Broyden--Fletcher--Goldfarb--Shanno (BFGS) optimization algorithm~\cite{nocedal2006numerical}. The forces $\mathbf{F}_i$ on each unconstrained atom $i$ are minimized until the maximum force component falls below a predefined threshold $f_{\text{max}}$:
\begin{equation}
    \max_{i} \| \mathbf{F}_i \|_{\infty} \le f_{\text{max}}
\end{equation}

During the optimization, selective constraints are applied. In the standard mode, the bottom fraction of the surface slab atoms (typically one-third of the $z$-axis dimension) is held fixed to simulate the bulk crystal termination, while the upper atomic layers and the adsorbate are fully relaxed.

Following physical simulation, the \textbf{Analyzer} systematically parses the relaxed trajectory to extract the precise local coordination environment.
As defined above, \textit{Chemical Slip} is detected when the relaxed structure exhibits a different surface-symbol fingerprint than the Planner's intent (see \S~2.2.1). The surface-symbol fingerprint is recomputed from the relaxed coordinates to ensure consistency. Canonical site-family mismatch is logged as a related geometric/site slip diagnostic rather than being conflated with the surface-symbol slip flag. 

The connectivity of the relaxed system is established by constructing a localized adjacency matrix $A$. Two atoms $i$ and $j$ are considered bonded if their interatomic distance $d_{ij}$ satisfies:
\begin{equation}
    d_{ij} \le \gamma \left( R_i + R_j \right)
\end{equation}
where $R_i$ and $R_j$ are the covalent radii of the atoms. For intramolecular dissociation checks, $\gamma=1.35$ is used to construct the adsorbate subgraph; for adsorbate--surface coordination, $\gamma=1.30$ is used by default and relaxed to $\gamma=1.45$ for strongly adsorbed states ($E_{\text{ads}} < -0.5$~eV). Using the adjacency matrix $A$, the system evaluates:
\begin{itemize}
    \item \textbf{Dissociation:} Computed via the connected components of the adsorbate subgraph. If the number of components exceeds one, intramolecular dissociation is flagged.
    \item \textbf{Coordination Environment:} By identifying all surface atoms $k$ bonded to the targeted adsorbate anchor atom $a$ ($A_{ak} = 1$), the actual site type is classified into \textit{ontop} ($\sum_k A_{ak} = 1$), \textit{bridge} ($\sum_k A_{ak} = 2$), or \textit{hollow} ($\sum_k A_{ak} \ge 3$).
\end{itemize}
If the realized site type or surface-symbol fingerprint differs from the \textbf{Planner}'s intent, the slip diagnostic is logged.

The \textbf{Analyzer} then updates the running memory with the adsorption energy, bond integrity, and slip status, enabling dynamic adjustment of the exploration strategy.

\subsection{Physical Backend: MACE-MP-0 Calculator}
\label{sec:mace_config}

The physical engine uses the MACE equivariant message-passing architecture and the MACE-MP foundation model~\cite{batatia2022mace,batatia2023foundation} as the surrogate calculator. Two configurations are employed:

\begin{itemize}
    \item \textbf{Primary benchmark}: MACE-MP-0 small checkpoint, CPU execution, float32 precision, dispersion disabled. This conservative setting reflects practical hardware constraints and is held fixed across all reported ablations and backend comparisons.
    \item \textbf{Sensitivity checks}: MACE-MP-0 large checkpoint, CUDA execution, float64 precision, dispersion enabled. These runs probe force-field sensitivity and are treated separately from the primary benchmark.
\end{itemize}

All primary benchmark and matched-baseline experiments use the same MACE-MP-0 small physical backend so that observed differences reflect search strategy, not energy-function choice. Key parameters and default settings are summarized in Table~\ref{tab:mace_config}.

\subsection{Evaluation Protocol and Experimental Setup}
\label{sec:eval_setup}

To isolate the contribution of individual mechanisms in AdsMind's closed-loop strategy, we design an ablation study with five variants:

\begin{itemize}
    \item \textbf{Full}: All mechanisms enabled (Slip feedback, Forbid constraint, Termination).
    \item \textbf{w/o Slip}: Removes slip-conditioned guidance and the associated negative constraint.
    \item \textbf{w/o Forbid}: Keeps slip detection but suppresses the explicit prohibition against retrying a failed site family.
    \item \textbf{w/o Term}: Disables convergence-based early stopping while preserving the rest of the loop.
    \item \textbf{1-Shot}: Disables all three mechanisms and sets the maximum attempt budget to one.
\end{itemize}

All agent-side benchmark runs in this manuscript use frozen experiment configurations rather than mutable product defaults. Four reasoning backends are evaluated: Gemini~2.5~Pro (\texttt{gemini-2.5-pro})~\cite{gemini2026modelcard}, Grok-4 (\texttt{grok-4-0709})~\cite{grok2025blog}, GPT-5.4 (\texttt{gpt-5.4-2026-03-05})~\cite{openai2026gpt5}, and Claude~Sonnet~4.6 (\texttt{claude-sonnet-4-6})~\cite{anthropic2026claude}, all run at zero temperature under fixed backend settings. For the primary benchmark protocol, the physical backend is locked to the MACE-MP-0 small configuration described in~\S~\ref{sec:mace_config}. Each frozen configuration records the model identifier and experiment settings needed to reproduce the run.

\subsection{CMU20 Experimental Setup}

The CMU20 benchmark uses the completed 20-case manifest, spanning atomic adsorbates (H), small multi-atom intermediates (OH and NNH), and larger intermediates (CH$_2$CH$_2$OH, OCHCH$_3$, and ONN(CH$_3$)$_2$) across intermetallic and monometallic surfaces.
The current CMU analysis uses the full 20-case manifest. Full runs use a maximum of five attempts. The 1-shot baseline uses the same tooling and physical backend but restricts the planner to a single attempt, allowing us to isolate the value of closed-loop refinement from the value of the base planner alone. Experiments are launched from frozen configurations with fixed seeds and structured result serialization for reproducibility. The five-variant ablation matrix (Full, w/o Slip, w/o Forbid, w/o Term, 1-Shot) is executed independently on all four LLM backends, yielding 400 total attempts (20 cases $\times$ 5 variants $\times$ 4 backends).
Natural failures (dissociation or reaction outcomes) are retained in the success-rate denominator; only external service or calculator errors are excluded from the analysis.

\subsection{OCD-GMAE62 Benchmark}

For validation on broader chemistry beyond CMU20, we assemble the \textbf{OCD-GMAE62} benchmark (\textbf{O}pen \textbf{C}hemistry \textbf{D}ataset -- \textbf{G}eneralized \textbf{M}achine-learning \textbf{A}dsorption \textbf{E}nergy 62-case dataset).
The benchmark contains 62 unique adsorption tasks spanning 48 intermetallic, 13 compound, and 1 monometallic surfaces, with 18 adsorbate labels including H, OH, CO, NO, N, \texttt{[N]=N}, \texttt{[N]=O}, \texttt{[NH]}, \texttt{[NH2]}, \texttt{[CH]=O}, \texttt{[CH2][O]}, \texttt{C[O]}, \texttt{C(=O)C}, and related carbon/nitrogen/oxygen intermediates.
The dataset is curated as a single 62-case benchmark; duplicated historical sampling records are not treated as additional benchmark cases.

The primary OCD-GMAE62 evaluation uses the same five-variant ablation matrix as CMU20 on all 62 cases and all four LLM backends.
This yields 1,240 raw attempts (62 cases $\times$ 5 variants $\times$ 4 backends) under the locked MACE-MP-0 small CPU protocol.
Metrics include success rate, success-conditioned $\Delta E = E_{\mathrm{variant}} - E_{\mathrm{Full}}$, four-backend range, reach-Full-within-0.01~eV, Wilcoxon signed-rank tests with Benjamini--Hochberg correction, and matched random/heuristic baseline comparisons.

A separate reproducibility analysis uses the 12 duplicated OCD-GMAE62 source records that appeared in both historical sampling batches.
These 12 cases are retained only as an overlap set for repeated-run analysis, not as a separate benchmark subset.
For each overlap case, we compare repeated runs across the same four LLM backends and five runtime variants, giving 240 repeated-run comparisons per additional repeat.
This design turns the historical overlap into a reproducibility audit while keeping the paper-facing benchmark definition as CMU20 plus OCD-GMAE62.

All experiments use the four LLM backends listed in Section~\ref{sec:eval_setup} under a fixed MACE-MP-0 small CPU protocol.
External service failures are excluded at the raw-attempt audit stage; all reported failures are natural dissociation or reaction outcomes.
